\documentclass[10pt,a4paper]{amsart}
\usepackage[margin=19mm]{geometry}
\usepackage{amsmath,amssymb,mathtools}
\usepackage[scr=rsfs,cal=euler]{mathalfa}
\usepackage{xfrac}
\usepackage[unicode,hidelinks,bookmarks=false]{hyperref}
\newcommand{\ie}{i.e.\ }
\newcommand{\cf}{cf.\ }
\newcommand{\resp}{respectively\ }
\newcommand{\Subsection}{\subsection}
\providecommand{\nobreakdash}{\nobreak\hbox{-}}
\setlength{\emergencystretch}{2em}
\multlinegap0pt
\usepackage{amssymb}
\usepackage{enumerate}
\newtheorem{theorem}{Theorem}
\newcommand{\be}{\mathbb E}
\newcommand{\ot}{\otimes}
\newcommand{\raro}{\rightarrow}
\newcommand{\wT}{\wt{T}}
\newcommand{\wt}{\widetilde}
\title{Completely contractive covariant representations of product system over $\mathbb N^2_0$}
\author{Dimple Saini}
\date{Source: arXiv:2606.22176v1 (CC BY 4.0)}
\begin{document}
\maketitle

\noindent\emph{Source and licence.} Completely contractive covariant representations of product system over $\mathbb N^2_0$. Version arXiv:2606.22176v1 (CC BY 4.0), distributed by arXiv under the Creative Commons Attribution 4.0 International licence, http://creativecommons.org/licenses/by/4.0/, as recorded in the arXiv licence metadata for this exact version and shown on the abstract page https://arxiv.org/abs/2606.22176v1. Full licence notice: \texttt{licenses/arxiv-2606.22176-CC-BY-4.0.txt}.\par\smallskip

\noindent\textit{Editorial scope.} Counted unit: the article's theorem JJ3 (source0.tex lines 313--315). The article's complete original proof of it is delivered with it (source0.tex lines 316--376, one \texttt{proof} environment of 61 lines). No mathematics is written, repaired or re-derived: the statement and the proof are the article's own lines, in the article's own notation, and no retained line is substituted or re-flowed.  Retained uncounted as the source-local context the proof needs: lines 183--189; lines 244--247; lines 271--276.  Nothing was omitted from any retained span, so the package carries no omission record.  No retained line contains a citation, so the wrapper carries no bibliography and no bibliography entry.  No retained line contains a figure environment, an image-inclusion command or drawing code, so no image is delivered and the package declares no asset dependency.  Editorial formatting only, in addition: an AMS \texttt{amsart} preamble that restates the article's own declarations and macros used by the retained lines, namely the environments \texttt{theorem} on separate counters, together with the standard packages those lines need.  The retained environments print in this document as Theorem 1 in order of appearance, whereas the article prints the same items with the numbering of its own full document.  The complete native source of the article as distributed in the arXiv e-print for this exact version is bundled unchanged as \texttt{sources/203329/source0.tex}.
\begin{theorem}\label{dilation}
		Assume $(\sigma, T)$ to be a c.c.c. representation of $E$ on $\mathcal{H},$ then the poisson kernel $\Pi$ is a contraction and
		\begin{align} \label{sid}
		  (I_{E} \ot \Pi)\widetilde{T}^{*}=\widetilde{S}^{*}\Pi  \:\: \mbox{and} \:\:  	\Pi \sigma(a)=\rho(a)\Pi, \:\:\:a \in \mathcal{B}. 
		\end{align}
		Moreover, $\Pi$ is an isometry if and only if $(\sigma, T) $ is pure. The induced representation $(\rho,S)$ of $E$ induced by $\sigma|_{\mathcal{D}_{T}}$ is called {\rm isometric dilation} of $(\sigma, T).$
\end{theorem}

	\begin{align}\label{BB1}
	&I_{\mathcal H}-\widetilde{T}^{(1)}\widetilde{T}^{(1)^*}+\widetilde{T}^{(1)}(I_{E_1}\otimes(I_{\mathcal{H}}-\widetilde{T}^{(2)}\widetilde{T}^{(2)^*}))\widetilde{T}^{(1)^*}=	I_{\mathcal H}-{\widetilde T}{\widetilde T^*}\\& \nonumber=I_{\mathcal H}-\widetilde{T}^{(2)}\widetilde{T}^{(2)^*}+
	\quad \widetilde{T}^{(2)}(I_{E_2}\otimes(I_{\mathcal{H}}-\widetilde{T}^{(1)}\widetilde{T}^{(1)^*}))\widetilde{T}^{(2)^*}.
	\end{align}

\begin{theorem}\label{identity}
	Using the above notations, let \begin{equation}
			\label{unitary} U = \begin{bmatrix}A&B\\C&0\end{bmatrix}:
			\mathcal{K} \oplus  (E_1 \ot\mathcal{D}_{ T^{(2)}})\to (E_2 \ot \mathcal{K}) \oplus  \mathcal{D}_{ T^{(2)}},
		\end{equation} be a unitary module map such that $$U(V\Delta(T^{(1)})h,(I_{E_1} \ot \Delta(T^{(2)}))\wT^{(1)^*}h)=((I_{E_2} \ot V\Delta(T^{(1)}))\wT^{(2)^*}h,\Delta(T^{(2)})h)$$ for all $h\in \mathcal{H}.$ Then there exists a contractive linear map $\widetilde{\Phi}:E_2\ot \mathcal{K}\to \mathcal{F}(E_1)\ot \mathcal{K}$ such that $$ (I_{E_2}\ot \Pi_V)\wT^{(2)^*}=\widetilde{M}_{\Phi}^*\Pi_V,$$ where $\widetilde{\Phi}=A^*+(I_{E_1}\ot C^*)B^*.$
	  \end{theorem}

	\begin{theorem}\label{JJ3}
	Assume $(\sigma, T^{(1)}, T^{(2)})$ to be a pure c.c.c. representation  of $\be$ over $\mathbb N^2_0$ on $\mathcal H$ such that $\mbox{dim~} (\mathcal{D}_{T^{(1)}} \oplus (E_1 \ot\mathcal{D}_{ T^{(2)}}))=\mbox{dim~} (\mathcal{D}_{ T^{(2)}}\oplus (E_2\ot \mathcal{D}_{ T^{(1)}}))< \infty.$ Then $(\sigma, T^{(1)}, T^{(2)})$ dilates to a pure isometric covariant representation.
\end{theorem}

\begin{proof}
	Since $(\sigma, T)$ is pure c.c.c. representation of $E:=E_1 \ot E_2$ on $\mathcal H,$ where $\widetilde{T}=\widetilde{T}^{(1)}(I_{E_1} \otimes \widetilde{T}^{(2)}).$ Let $\Pi: \mathcal{H} \longrightarrow \mathcal{F}(E)\otimes\mathcal{D}_{T}$ be the isometric dilation of $(\sigma,T)$ (see Theorem \ref{dilation}). Let $\mathcal K= \mathcal{D}_{T^{(2)}} \oplus(E_2 \ot\mathcal{D}_{T^{(1)}}),$ from Equation (\ref{BB1}), we define an isometric module map $V:\mathcal{D}_{T}\to  \mathcal K$ by
	$$V(\Delta(T)h):=(\Delta(T^{(2)})h,  (I_{E_2} \ot \Delta(T^{(1)}))\wT^{(2)^*}h)$$ for $h\in \mathcal{H}.$ Consequently, $\Pi_V:=(I_{\mathcal{F}(E)}\ot V)\Pi: \mathcal{H} \longrightarrow \mathcal{F}(E)\otimes\mathcal{K}$ is an isometric dilation of $(\sigma,T).$ Thus $\wT\cong P_{\mathcal{Q}}\widetilde{S}|_{E\ot \mathcal{Q}}$, where $\mathcal{Q}=\Pi_V\mathcal{H},$ $(\rho,S)$ is an induced representation of $E$ and $\widetilde{S}^{*}(\mathcal{Q})\subseteq E\ot \mathcal{Q}$ (see Theorem \ref{identity}). 
	Since $$(I_{E_2}\ot V\Delta(T))\wT^{(2)^*}h=((I_{E_2}\ot \Delta(T^{(2)}))\wT^{(2)^*}h,(I_{E_2^{\ot 2}}\ot \Delta(T^{(1)}))\wT_2^{(2)^*}h),\quad  h\in \mathcal{H}.$$
	Let $j_1 : \mathcal{D}_{T^{(2)}} \raro \mathcal{K}$ and $j_2 :  E_2 \ot\mathcal{D}_{ T^{(1)}} \raro \mathcal{K}$ be
	the inclusion maps, defined by
	\[
	j_1(h_1) := (h_1,0)\quad  \mbox{and} \quad j_2 (h_2) := (0,h_2)
	\quad (h_1 \in \mathcal{D}_{T^{(2)}}, h_2 \in  E_2 \ot\mathcal{D}_{ T^{(1)}}).
	\]
	Then $P : = j_1 j_1^* \in B(\mathcal{K})$ is an orthogonal
	projection of $\mathcal{K}$ onto $ \mathcal{D}_{ T^{(2)}}$, i.e.,
	\[
	P(h_1,h_2) = (h_1,0) \quad \quad \quad (h_1,h_2) \in \mathcal{K}.
	\]
	Thus, $j_2 j_2^* = P^\perp$ is the orthogonal projection of $\mathcal{K}$
	onto $E_2 \ot\mathcal{D}_{ T^{(1)}}.$ Since $\mbox{dim~} (\mathcal{D}_{T^{(1)}} \oplus (E_1 \ot\mathcal{D}_{ T^{(2)}}))< \infty$ and $\mbox{dim~} (\mathcal{D}_{ T^{(2)}}\oplus (E_2\ot \mathcal{D}_{ T^{(1)}}))< \infty.$ Define a
	unitary module map $U:E_2\ot\mathcal{K}\to E_2\ot E_1 \ot \mathcal{D}_{ T^{(2)}}\oplus E_2\ot \mathcal{D}_{ T^{(1)}}$ by
	\begin{align*}U((I_{E_2}\ot \Delta(T^{(2)}))\wT^{(2)^*}h,(I_{E_2^{\ot 2}} \ot \Delta(T^{(1)}))\wT_2^{(2)^*}h):=((t_{1,2}\ot \Delta(T^{(2)}))\wT^{*}h,(I_{E_2}\ot \Delta(T^{(1)}))\wT^{(2)^*}h).\end{align*} 
	Then
	\[
	U_1=
	\begin{bmatrix}
		U^*P^{\perp}& U^*(t_{1,2}\ot j_1)\\
		j_1^*& 0
	\end{bmatrix}: \mathcal{K} \oplus  (E \ot\mathcal{D}_{ T^{(2)}})\to (E_2 \ot \mathcal{K}) \oplus  \mathcal{D}_{ T^{(2)}}
	\] is a unitary module map. For $h\in \mathcal{H}$ we have
	\begin{align*}
		&U_1(V\Delta(T)h,(I_{E}\ot \Delta(T^{(2)}))\wT^{*}h)=U_1(\Delta(T^{(2)})h,  (I_{E_2} \ot \Delta(T^{(1)}))\wT^{(2)^*}h,(I_{E}\ot \Delta(T^{(2)}))\wT^{*}h)\\&=(U^*((t_{1,2}\ot \Delta(T^{(2)}))\wT^{*}h,(I_{E_2}\ot \Delta(T^{(1)}))\wT^{(2)^*}h),\Delta(T^{(2)})h)\\&=((I_{E_2}\ot \Delta(T^{(2)}))\wT^{(2)^*}h,(I_{E_2^{\ot 2}} \ot \Delta(T^{(1)}))\wT_2^{(2)^*}h,\Delta(T^{(2)})h)\\&=((I_{E_2}\ot V\Delta(T))\wT^{(2)^*}h,\Delta(T^{(2)})h).
	\end{align*} From Theorem \ref{identity}, there exists a contractive linear map $\widetilde{\Phi}_2:E_2 \ot \mathcal{K}\to \mathcal{F}(E)\ot \mathcal{K}$ such that $$(I_{E_2}\ot \Pi_V)\wT^{(2)^*}=\widetilde{M}_{{\Phi}_2}^*\Pi_V,$$ where $\widetilde{\Phi}_2=P^{\perp}U+(I_{E}\ot j_1)(t_{2,1}\otimes j_1^*)U=(P^{\perp} + \widetilde{S}(t_{2,1}\otimes P))U.$ 

Similarly, since $\mathcal{D}_{T^{(1)}} \oplus (E_1 \ot\mathcal{D}_{ T^{(2)}})\cong \mathcal{D}_{ T^{(2)}}\oplus (E_2\ot \mathcal{D}_{ T^{(1)}}),$  let $\mathcal{K}=\mathcal{D}_{T^{(1)}} \oplus(E_1 \ot\mathcal{D}_{T^{(2)}})$ and 
$V(\Delta(T)h):=(\Delta(T^{(1)})h,  (I_{E_1} \ot \Delta(T^{(2)}))\wT^{(1)^*}h)$ for $h\in \mathcal{H}.$ 
Since $$(I_{E_1}\ot V\Delta(T))\wT^{(1)^*}h=((I_{E_1}\ot \Delta(T^{(1)}))\wT^{(1)^*}h,(I_{E_1^{\ot 2}}\ot \Delta(T^{(2)}))\wT_2^{(1)^*}h), \quad  h\in \mathcal{H}.$$
Let $i_1 : \mathcal{D}_{T^{(1)}} \raro \mathcal{K}$ and $i_2 :  E_1 \ot\mathcal{D}_{ T^{(2)}} \raro \mathcal{K}$ be
the inclusion maps, defined by
\[
i_1(h_1) := (h_1,0)\quad  \mbox{and} \quad i_2 (h_2) := (0,h_2)
\quad (h_1 \in \mathcal{D}_{T^{(1)}}, h_2 \in  E_1 \ot\mathcal{D}_{T^{(2)}}).
\]
Then $Q : = i_1 i_1^* \in B(\mathcal{K})$ is an orthogonal
projection of $\mathcal{K}$ onto $ \mathcal{D}_{ T^{(1)}}$, i.e.,
\[
Q(h_1,h_2) = (h_1,0) \quad \quad \quad (h_1,h_2) \in \mathcal{K}.
\]
Thus $i_2 i_2^* = Q^\perp$ is the orthogonal projection of $\mathcal{K}$
onto $E_1 \ot\mathcal{D}_{ T^{(2)}}.$ Define a
unitary module map $U':E_1\ot\mathcal{K}\to E_1\ot E_2 \ot \mathcal{D}_{ T^{(1)}}\oplus E_1\ot \mathcal{D}_{ T^{(2)}}$ by
\begin{align*}U'((I_{E_1}\ot \Delta(T^{(1)}))\wT^{(1)^*}h,(I_{E_1^{\ot 2}} \ot \Delta(T^{(2)}))\wT_2^{(1)^*}h):=((I_E\ot \Delta(T^{(1)}))\wT^{*}h,(I_{E_1}\ot \Delta(T^{(2)}))\wT^{(1)^*}h).\end{align*} 
Then
\[
U_2=
\begin{bmatrix}
	U'^*Q^{\perp}& U'^*(I_E\ot i_1)\\
	i_1^*& 0
\end{bmatrix}: \mathcal{K} \oplus  (E \ot\mathcal{D}_{ T^{(1)}})\to (E_1 \ot \mathcal{K}) \oplus  \mathcal{D}_{ T^{(1)}}
\] is a unitary module map. For $h\in \mathcal{H}$ we have
\begin{align*}
	&U_2(V\Delta(T)h,(I_{E}\ot \Delta(T^{(1)}))\wT^{*}h)=U_2(\Delta(T^{(1)})h,  (I_{E_1} \ot \Delta(T^{(2)}))\wT^{(1)^*}h,(I_{E}\ot \Delta(T^{(1)}))\wT^{*}h)\\&=(U'^*((I_E\ot \Delta(T^{(1)}))\wT^{*}h,(I_{E_1}\ot \Delta(T^{(2)}))\wT^{(1)^*}h),\Delta(T^{(1)})h)\\&=((I_{E_1}\ot \Delta(T^{(1)}))\wT^{(1)^*}h,(I_{E_1^{\ot 2}} \ot \Delta(T^{(2)}))\wT_2^{(1)^*}h,\Delta(T^{(1)})h)\\&=((I_{E_1}\ot V\Delta(T))\wT^{(1)^*}h,\Delta(T^{(1)})h).
\end{align*} From Theorem \ref{identity}, there exists a contractive linear map $\widetilde{\Phi}_1:E_1 \ot \mathcal{K}\to \mathcal{F}(E)\ot \mathcal{K}$ such that $$(I_{E_1}\ot \Pi_{V})\wT^{(1)^*}=\widetilde{M}_{{\Phi}_1}^*\Pi_{V},$$ where $\widetilde{\Phi}_1=Q^{\perp}U'+(I_{E}\ot i_1i_1^*)U'=(Q^{\perp} + \widetilde{S}(I_E\otimes Q))U'.$ We conclude that $(\rho,M_{\Phi_1},M_{\Phi_2})$ is the pure isometric dilation of $(\sigma,T^{(1)},T^{(2)}).$ This complets the proof.
\end{proof}

\end{document}
